\documentclass[conference]{IEEEtran}

\usepackage[utf8]{inputenc}
\usepackage[T1]{fontenc}
\usepackage[spanish,es-tabla]{babel}
\usepackage{cite}
\usepackage{amsmath}
\usepackage{booktabs}
\usepackage{array}
\usepackage{url}
\usepackage{graphicx}
\usepackage{tikz}
\usetikzlibrary{arrows.meta,positioning}

\title{Implementación y Análisis de Inferencia de E2D2: un Transformer Encoder--Decoder de Difusión Discreta para Traducción Automática}

\author{
\IEEEauthorblockA{
Procesamiento de Datos Secuenciales con Deep Learning\\
Universidad Autónoma Occidente\\
Especialización 202601
}\\[1ex]
\begin{tabular}{>{\centering\arraybackslash}p{0.22\textwidth}
                >{\centering\arraybackslash}p{0.22\textwidth}
                >{\centering\arraybackslash}p{0.22\textwidth}
                >{\centering\arraybackslash}p{0.22\textwidth}}
Nicolás Vásquez Renjifo &
Jorge Luis Fong Gutierrez &
Jhonatan David Rengifo Bermeo &
Mateo González Ruiz \\
\texttt{\scriptsize nicolas.vasquez@uao.edu.co} &
\texttt{\scriptsize jorge.fong@uao.edu.co} &
\texttt{\scriptsize jhonatan.rengifo@uao.edu.co} &
\texttt{\scriptsize mateo.gonzalez@uao.edu.co} \\
Ingeniero Mecatrónico &
Ingeniero Mecatrónico &
Ingeniero Mecatrónico &
Ingeniero Mecatrónico
\end{tabular}
}

\begin{document}

\maketitle

\begin{abstract}
Este trabajo presenta la implementación y evaluación experimental de E2D2, una arquitectura Transformer encoder--decoder basada en modelos de difusión discreta para lenguaje natural. El proyecto responde a la consigna de implementar inferencia sobre una arquitectura Transformer encoder--decoder aplicada a datos secuenciales complejos, sin entrenar desde cero, utilizando un artículo científico con código y pesos preentrenados disponibles. Se utiliza el checkpoint \texttt{kuleshov-group/e2d2-wmt}, entrenado sobre WMT14 de-en, y se evalúa su comportamiento en traducción alemán--inglés. El estudio incluye inspección interna de la arquitectura, inferencia en Google Colab, un estudio de ablación sobre \texttt{block\_size}, pasos de difusión y caché, comparación extendida con el baseline autoregresivo \texttt{Helsinki-NLP/opus-mt-de-en}, y análisis cualitativo de errores. Los resultados muestran que la configuración oficial de E2D2, \texttt{block\_size=4} y \texttt{num\_steps=4}, es la más equilibrada entre las variantes E2D2 evaluadas. En la evaluación extendida de 150 ejemplos, E2D2 obtiene BLEU 16.53 y chrF 49.66, mientras que OPUS-MT obtiene BLEU 25.02 y chrF 54.09. Estos resultados permiten discutir críticamente las diferencias entre decodificación autoregresiva y generación por denoising discreto.
\end{abstract}

\begin{IEEEkeywords}
Transformer, encoder--decoder, difusión discreta, E2D2, traducción automática, inferencia, WMT14.
\end{IEEEkeywords}

\section{Introducción}

Los modelos Transformer han dominado gran parte del procesamiento de lenguaje natural debido a su capacidad para modelar dependencias de largo alcance en secuencias. Tradicionalmente, las tareas de generación de texto, como traducción o resumen, se han abordado con modelos autoregresivos que generan un token a la vez. Aunque este paradigma es efectivo, su naturaleza secuencial limita la paralelización durante la inferencia.

El artículo \emph{Encoder-Decoder Diffusion Language Models for Efficient Training and Inference} propone E2D2, una arquitectura Transformer encoder--decoder que adapta modelos de difusión discreta al modelado de lenguaje \cite{arriola2025e2d2}. La idea central es separar dos tipos de computación: la representación de tokens limpios mediante un encoder grande, y el denoising iterativo de tokens corruptos mediante un decoder más liviano. Esta separación busca mejorar la eficiencia de inferencia respecto a modelos de difusión decoder-only.

Este proyecto implementa inferencia con E2D2 sobre traducción alemán--inglés. La tarea se considera adecuada para la consigna porque el lenguaje natural es una secuencia compleja y no corresponde a series de tiempo tabulares o sensores. Además, el paper seleccionado cuenta con repositorio oficial, integración en Hugging Face y pesos preentrenados disponibles \cite{kuleshov_e2d2_repo,hf_e2d2_wmt}.

\section{Objetivos}

El objetivo general es implementar y analizar una arquitectura Transformer encoder--decoder aplicada a datos secuenciales complejos, usando inferencia con pesos preentrenados.

Los objetivos específicos son:

\begin{itemize}
    \item Revisar la arquitectura E2D2 propuesta en el artículo base.
    \item Cargar pesos preentrenados oficiales y ejecutar inferencia en Google Colab.
    \item Inspeccionar la configuración interna del modelo cargado.
    \item Evaluar el impacto de \texttt{block\_size}, \texttt{num\_steps} y uso de caché.
    \item Comparar E2D2 contra un baseline Transformer encoder--decoder autoregresivo.
    \item Analizar errores cualitativos de traducción y discutir limitaciones.
\end{itemize}

\section{Cumplimiento de la Consigna}

La Tabla \ref{tab:consigna} resume cómo el proyecto satisface explícitamente los requisitos de la consigna.

\begin{table*}[htbp]
\caption{Verificación de cumplimiento de la consigna}
\label{tab:consigna}
\centering
\begin{tabular}{p{0.28\textwidth}p{0.62\textwidth}}
\toprule
\textbf{Requisito} & \textbf{Cumplimiento en el proyecto} \\
\midrule
Arquitectura Transformer encoder--decoder & Se implementa inferencia con E2D2, una arquitectura con encoder Transformer y decoder Transformer de difusión. \\
Datos secuenciales complejos & Se trabaja con lenguaje natural, específicamente traducción automática alemán--inglés. \\
No usar series de tiempo tradicionales & No se usan datos tabulares, sensores ni registros estructurados. \\
Artículo científico base & El proyecto se basa en el paper \emph{Encoder-Decoder Diffusion Language Models for Efficient Training and Inference}. \\
Código disponible & Se usa el repositorio oficial publicado por Kuleshov Group. \\
Pesos preentrenados & Se usa el checkpoint público \texttt{kuleshov-group/e2d2-wmt}. \\
No entrenar desde cero & Solo se ejecuta inferencia y evaluación con pesos preentrenados. \\
Explicación profunda & Se documentan arquitectura, block diffusion, máscaras conceptuales, métricas, ablación, baseline y errores. \\
\bottomrule
\end{tabular}
\end{table*}

\section{Arquitectura E2D2}

E2D2 es una arquitectura encoder--decoder diseñada para modelos de difusión discreta en lenguaje. A diferencia de un Transformer encoder--decoder clásico, donde el decoder suele generar tokens autoregresivamente, E2D2 utiliza un decoder de difusión que refina tokens ruidosos o enmascarados.

\subsection{Componentes}

La arquitectura divide el procesamiento en dos módulos principales:

\begin{itemize}
    \item \textbf{Encoder:} procesa tokens limpios, como el prompt, la fuente alemana o bloques ya generados. En el checkpoint WMT usado, el encoder tiene 28 capas.
    \item \textbf{Decoder:} realiza denoising sobre el bloque activo. En el checkpoint WMT usado, el decoder tiene 4 capas.
\end{itemize}

Esta asimetría permite que el encoder, más costoso, se invoque con menor frecuencia, mientras que el decoder, más pequeño, puede ejecutarse múltiples veces durante el refinamiento de cada bloque.

\begin{figure}[htbp]
\centering
\begin{tikzpicture}[
    node distance=0.28cm,
    box/.style={draw, rounded corners, align=center, minimum width=3.15cm, minimum height=0.75cm, font=\footnotesize},
    arrow/.style={-{Latex[length=2mm]}, thick}
]
\node[box] (src) {Texto fuente\\en alemán};
\node[box, below=of src] (enc) {Encoder Transformer\\grande: 28 capas};
\node[box, below=of enc] (h) {Representaciones\\limpias $h$};
\node[box, below=of h] (dec) {Decoder Transformer\\difusión: 4 capas};
\node[box, below=of dec] (block) {Bloque objetivo ruidoso\\$[MASK,\ldots,MASK]$};
\node[box, below=of block] (out) {Traducción\\inglés};
\draw[arrow] (src) -- (enc);
\draw[arrow] (enc) -- (h);
\draw[arrow] (h) -- (dec);
\draw[arrow] (block) -- (dec);
\draw[arrow] (dec) -- (out);
\draw[arrow] (dec.east) .. controls +(1.0,0.0) and +(1.0,0.0) .. node[right, font=\scriptsize, align=center]{T pasos\\de denoising} (block.east);
\end{tikzpicture}
\caption{Flujo conceptual de inferencia en E2D2. El encoder representa tokens limpios y el decoder refina tokens ruidosos por bloques.}
\label{fig:e2d2_diagrama}
\end{figure}

\subsection{Block Diffusion}

En block diffusion, la secuencia se divide en bloques de tamaño $S$. Para cada bloque, el modelo parte de una representación ruidosa y aplica pasos de denoising. La inferencia puede representarse conceptualmente como:

\begin{verbatim}
1. Codificar la secuencia fuente con el encoder.
2. Inicializar un bloque objetivo con tokens ruidosos.
3. Ejecutar T pasos de denoising con el decoder.
4. Agregar el bloque generado al contexto.
5. Repetir hasta completar la secuencia.
\end{verbatim}

El checkpoint E2D2-WMT utiliza \texttt{block\_size=4} y \texttt{num\_steps=4} como configuración oficial.

\section{Metodología}

\subsection{Entorno}

Los experimentos se ejecutaron en Google Colab con GPU Tesla T4 y 15.64 GB de VRAM. Aunque el notebook incluye detección experimental de TPU/XLA, la corrida registrada se realizó con CUDA. Las dependencias principales fueron \texttt{transformers==4.52.4}, \texttt{accelerate==1.11.0}, \texttt{datasets==4.2.0}, \texttt{evaluate}, \texttt{sacrebleu} y \texttt{sentencepiece}.

\subsection{Modelo Principal}

El modelo principal fue \texttt{kuleshov-group/e2d2-wmt}, publicado en Hugging Face. La Tabla \ref{tab:arquitectura} resume los atributos inspeccionados directamente desde el checkpoint.

Para reproducibilidad, se fijó la revisión exacta del checkpoint E2D2-WMT en Hugging Face: \texttt{7c65309acc0d6d89049c9aca4ea02f9d10d1da5a}. El baseline OPUS-MT también se fijó a la revisión \texttt{1a922f3b32a8e809e17a47d4b32142d8105924e5}. Esto evita que una actualización remota altere silenciosamente los resultados.

\begin{table}[htbp]
\caption{Configuración inspeccionada del checkpoint E2D2-WMT}
\label{tab:arquitectura}
\centering
\begin{tabular}{ll}
\toprule
\textbf{Atributo} & \textbf{Valor} \\
\midrule
Tipo de modelo & E2D2 \\
Tipo de difusión & absorbing \\
Capas del encoder & 28 \\
Capas del decoder & 4 \\
Hidden size & 512 \\
Intermediate size & 1536 \\
Block size & 4 \\
Longitud máxima & 256 \\
Tokenizer & Qwen/Qwen3-0.6B-Base \\
Parámetros totales & 331.8M \\
\bottomrule
\end{tabular}
\end{table}

\subsection{Dataset}

Se utilizó una muestra inicial de 30 ejemplos del conjunto de prueba WMT14 de-en para el estudio de ablación y análisis cualitativo. Adicionalmente, el notebook final incluye una sección de evaluación extendida configurada para 150 ejemplos, enfocada en comparar la configuración oficial de E2D2 contra el baseline OPUS-MT. Cada ejemplo contiene una fuente en alemán y una referencia humana en inglés. El tamaño de muestra extendido busca mejorar la robustez empírica sin convertir el proyecto en una reproducción completa de todas las evaluaciones del paper.

\subsection{Baseline}

Para contrastar E2D2 con una arquitectura encoder--decoder autoregresiva tradicional, se utilizó \texttt{Helsinki-NLP/opus-mt-de-en}, un modelo OPUS-MT alemán--inglés \cite{tiedemann2020opusmt}. Este baseline tiene aproximadamente 74.4M parámetros y genera traducciones mediante decodificación autoregresiva.

\subsection{Métricas}

Se reportaron las siguientes métricas:

\begin{itemize}
    \item \textbf{BLEU:} métrica estándar de traducción automática basada en coincidencias de n-gramas \cite{post2018call}.
    \item \textbf{chrF:} métrica basada en F-score de caracteres, útil para variaciones morfológicas y lexicales \cite{popovic2015chrf}.
    \item \textbf{Tiempo promedio:} segundos por ejemplo.
    \item \textbf{Tokens por segundo:} throughput de generación.
    \item \textbf{Repetición de bigramas:} tasa simple de bigramas repetidos en la salida.
\end{itemize}

\section{Implementación}

La inferencia se implementó en un notebook de Colab. La entrada alemana se construyó como:

\begin{verbatim}
BOS + texto_aleman + EOS
\end{verbatim}

Luego, el modelo E2D2 generó la continuación en inglés. Para cada traducción se midió tiempo, número de tokens generados, tokens por segundo y texto resultante. También se exportaron archivos CSV con resultados por configuración y análisis cualitativo.

El notebook final documentado contiene comentarios en las líneas importantes del código: carga de dependencias, selección de dispositivo, carga del checkpoint, construcción del prompt, generación E2D2, métricas, ablación, baseline OPUS-MT y exportación de resultados.

\section{Resultados}

\subsection{Ablation Study de E2D2}

La Tabla \ref{tab:ablation} muestra los resultados de las configuraciones E2D2 evaluadas. La configuración oficial del checkpoint fue la más equilibrada.

\begin{table*}[htbp]
\caption{Ablation study de E2D2 sobre 30 ejemplos WMT14 de-en}
\label{tab:ablation}
\centering
\begin{tabular}{lrrrrrr}
\toprule
\textbf{Configuración} & \textbf{BLEU} & \textbf{chrF} & \textbf{Tiempo (s)} & \textbf{Tokens/s} & \textbf{Tokens gen.} & \textbf{Rep. bigrama} \\
\midrule
E2D2 b2/s2 cache & 0.00 & 0.31 & 0.140 & 15.75 & 2.00 & 0.033 \\
E2D2 b4/s4 cache & 17.85 & 50.00 & 1.066 & 37.24 & 36.53 & 0.092 \\
E2D2 b8/s8 cache & 10.13 & 45.65 & 0.950 & 71.03 & 65.87 & 0.240 \\
E2D2 b4/s4 no-cache & -- & -- & -- & -- & -- & -- \\
\bottomrule
\end{tabular}
\end{table*}

La configuración \texttt{b2/s2} generó salidas demasiado cortas, con apenas dos tokens en promedio. La configuración \texttt{b8/s8} aumentó los tokens por segundo, pero deterioró BLEU y elevó la repetición. La ruta \texttt{use\_cache=False} falló en todos los ejemplos con un error de dimensiones, por lo que se reporta como limitación técnica de esta implementación y no como resultado competitivo.

\subsection{Comparación Extendida con OPUS-MT}

La Tabla \ref{tab:comparacion} compara E2D2 con el baseline autoregresivo OPUS-MT sobre 150 ejemplos del conjunto de prueba WMT14 de-en. Esta comparación extendida se usa como resultado principal de evaluación, mientras que la Tabla \ref{tab:ablation} se mantiene como estudio de sensibilidad de parámetros.

\begin{table*}[htbp]
\caption{Comparación extendida entre E2D2 y baseline autoregresivo sobre 150 ejemplos}
\label{tab:comparacion}
\centering
\begin{tabular}{lrrrrr}
\toprule
\textbf{Modelo} & \textbf{BLEU} & \textbf{chrF} & \textbf{Tiempo (s)} & \textbf{Tokens/s} & \textbf{Rep. bigrama} \\
\midrule
E2D2 b4/s4 cache & 16.53 & 49.66 & 0.706 & 50.32 & 0.077 \\
OPUS-MT autoregresivo & 25.02 & 54.09 & 0.189 & 144.68 & 0.005 \\
\bottomrule
\end{tabular}
\end{table*}

En esta muestra extendida, OPUS-MT supera a E2D2 en calidad automática y velocidad. Esto no invalida E2D2 como objeto de estudio, ya que el propósito principal del proyecto es comprender e implementar una arquitectura encoder--decoder de difusión discreta con pesos preentrenados. La comparación permite discutir diferencias arquitectónicas: OPUS-MT es un traductor autoregresivo especializado, mientras E2D2 genera mediante refinamiento discreto por bloques.

\section{Análisis Cualitativo}

El análisis cualitativo muestra que E2D2 puede generar traducciones razonables, pero también presenta errores de repetición y continuaciones no justificadas. Por ejemplo, para la frase alemana:

\begin{quote}
\emph{Pro Fahrtrichtung gibt es drei Lichtanlagen.}
\end{quote}

La referencia es:

\begin{quote}
\emph{There are three sets of lights per direction of travel.}
\end{quote}

E2D2 produjo:

\begin{quote}
\emph{There are three light systems per direction. The car is 100 m long and 100 m long...}
\end{quote}

OPUS-MT produjo:

\begin{quote}
\emph{There are three lighting systems per direction of travel.}
\end{quote}

Este caso evidencia una limitación práctica del sampling E2D2 en esta configuración: el modelo puede traducir correctamente el inicio, pero luego extenderse con contenido repetitivo o ajeno a la fuente. Por ello se incluyó una métrica simple de repetición de bigramas.

\section{Discusión}

Los resultados confirman que la configuración oficial de E2D2-WMT es la más razonable dentro de las variantes evaluadas. La reducción de \texttt{block\_size} a 2 produce generación insuficiente, mientras que aumentarlo a 8 incrementa la longitud y la repetición. Esto es coherente con el trade-off reportado en el paper: el tamaño de bloque afecta calidad, velocidad y frecuencia con la que se actualiza el contexto.

La comparación con OPUS-MT muestra que un modelo autoregresivo especializado puede ser más fuerte en una evaluación acotada de traducción. Sin embargo, el interés de E2D2 no está en reemplazar directamente a todos los traductores autoregresivos, sino en explorar una arquitectura alternativa donde la generación se formula como denoising discreto. Esta perspectiva es relevante para estudiar nuevas formas de paralelización y eficiencia en modelos de lenguaje.

\section{Limitaciones}

Las principales limitaciones del proyecto son:

\begin{itemize}
    \item Las ablaciones se ejecutan sobre una muestra reducida para controlar el tiempo de cómputo en Colab; la comparación principal E2D2 vs. OPUS-MT se ejecutó sobre 150 ejemplos.
    \item No se reproduce el entrenamiento desde cero porque la consigna prioriza inferencia con pesos preentrenados y análisis arquitectónico; entrenar E2D2-WMT excede los recursos razonables de Colab.
    \item El baseline OPUS-MT es más pequeño pero altamente especializado para traducción.
    \item Las métricas automáticas en muestras acotadas pueden variar y deben complementarse con análisis cualitativo.
\end{itemize}

\section{Conclusiones}

El proyecto implementó exitosamente inferencia con E2D2, una arquitectura Transformer encoder--decoder de difusión discreta aplicada a traducción automática. Se verificó internamente la estructura del checkpoint, confirmando un encoder de 28 capas y un decoder de 4 capas, con 331.8M parámetros. La configuración oficial \texttt{block\_size=4}, \texttt{num\_steps=4} y caché activado fue la variante E2D2 más equilibrada. La comparación con OPUS-MT mostró que el baseline autoregresivo fue superior en esta muestra, pero también permitió contextualizar las fortalezas y debilidades de la generación por difusión discreta.

En conjunto, la entrega cumple la consigna: se basa en un paper científico, usa código y pesos preentrenados disponibles, no entrena desde cero, implementa inferencia, trabaja con datos secuenciales complejos y explica en profundidad el funcionamiento de una arquitectura Transformer encoder--decoder.

\bibliographystyle{IEEEtran}
\bibliography{Informe_E2D2_IEEE}

\end{document}
